\appendix
\section{提高专题与模型扩展}
本附录承接原资料中的提高内容。建议在完成正文例题后选择学习。每次变换模型时，先检查状态语义是否改变，不要仅凭代码形式相似就替换循环。

\subsection{多重背包与二进制拆分}
\practice{P1776}{宝物筛选}
每种物品有数量上限，拆成01物品包；输入字段次序以原题为准。

第i种物品最多有 $s_i$ 件。朴素转移枚举件数k：
\[f[i][v]=\max_{0\le k\le s_i,\ kc_i\le v}\{f[i-1][v-kc_i]+kw_i\}.\]
当数量较大时，把 $s_i$ 件拆成若干只能选一次的物品包。例如13拆成1、2、4、6，每包体积和价值乘相应件数。选择这些包可以表示0到13的任意件数，转化为01背包。拆包数量为 $O(\log(s_i+1))$，总时间为 $O(V\sum_i\log(s_i+1))$。

\begin{lstlisting}
int remaining = count;
for (int batch = 1; remaining > 0; batch *= 2) {
    int take = min(batch, remaining);
    remaining -= take;
    long long cc = 1LL * take * cost;
    long long ww = 1LL * take * value;
    if (cc <= capacity)
        for (int v = capacity; v >= cc; --v)
            f[v] = max(f[v], f[v-(int)cc] + ww);
}
\end{lstlisting}
此片段假设count不超过 $10^6$、体积为正，避免batch翻倍溢出。若采用更大数量范围，应使用更宽类型并检查翻倍边界。

\textbf{为什么能覆盖。}若已有包可以组成0到t，再加入大小不超过t+1的新包，就可以连续覆盖到t加新包大小。按1、2、4增长并保留余数满足这个条件。

\textbf{计数陷阱。}这个拆分通常用于最优值或可达性，不能直接拿来统计原问题的方案数。例如拆成1、2、4、6后，6可由6或2+4表示，形成不同拆包选择却代表同样件数。计数背包应另行设计不重不漏的转移。

\subsection{混合背包与多维费用}
\practice{P1507}{NASA的食物计划}
对应本节二维费用模型，两种资源容量都必须满足。

混合背包对每种物品分别判断：限一次用逆序，无限次用正序，有限次可枚举或拆包；组内互斥则按整组处理，不能把同组物品当普通独立物品。不同模型混在一起时，应确保每轮转移的“一个决策单位”仍然清晰。

若每件物品同时占体积c和重量d，容量限制为V和W，定义 $f[v][w]$ 为两种资源限制下的最大价值。对于01物品，两个容量维度都逆序：
\begin{lstlisting}
for (int v = V; v >= cost; --v)
    for (int w = W; w >= weight; --w)
        f[v][w] = max(f[v][w], f[v-cost][w-weight]+value);
\end{lstlisting}
时间 $O(nVW)$，空间 $O(VW)$。增加状态维度会迅速增加计算量。应选择取值范围较小的量作维度，必要时考虑按价值记录最小成本，但也要重新定义答案和不可达值。

\subsection{方案恢复}
\practice{P1064}{金明的预算方案}
这是依赖背包的拓展题：附件必须依附主件。先枚举每个主件的合法购买组合，再作分组背包；不是普通独立物品模型。

若要求输出选了哪些物品，仅得到最大值还不够。对于二维01背包，可在每个状态记录是否选择第i件，或保留整张f表，从 $(n,V)$ 逆推：若最优值来自不选则i减1；若来自选择，则记录i并把容量减去体积。

一维压缩后直接保存“当前容量的父指针”容易被后续物品覆盖，造成恢复结果重复选同一物品。为了初学者的可靠实现，先使用二维记录恢复，再考虑持久化等高级优化。并列最优解允许任选时，固定一种规则即可；要求字典序最小时必须把并列规则也纳入设计。

\subsection{计数背包的撤销}
\practice{P4141}{消失之物}
删除的是一个指定物品，输出恰好装满的方案数末位数字；使用本节计数撤销公式。

设F[s]是所有物品组成体积s的01选择方案数，删除一个体积为w的\textbf{指定物品}后，G[s]是剩余方案数。按是否选该物品分类：
\[F[s]=G[s]+G[s-w],\qquad G[s]=F[s]-G[s-w].\]
负下标项为0。因G[s]依赖更小的G，按s从小到大计算。若只输出末位数字，应在模10下减法并把负数归一化：
\begin{lstlisting}
for (int s = 0; s <= capacity; s++) g[s] = all[s];
for (int s = w; s <= capacity; ++s)
    g[s] = (all[s] - g[s-w] + 10) % 10;
\end{lstlisting}
例如两个体积都为1的不同物品，完整方案数为 $F[0..2]=[1,2,1]$；删去其中一件，得到 $G=[1,1,0]$。物品即使体积相同也按身份区分。该线性计数关系可撤销，但最大值背包的max会丢失候选信息，不能简单相减恢复。

\subsection{LIS优化与排列LCS}
\practice{P1439}{两个排列的最长公共子序列}
利用元素各出现一次的性质映射成LIS；二次LCS不能覆盖大规模数据。

定义tail[len]为所有长度len的严格上升子序列中最小的末尾值。这个数组严格递增。处理x时，用lower\_bound找到第一个不小于x的位置并替换；若不存在就把x接到末尾。替换更小的末尾不会改变已获得的长度，却给后续元素留下更多延伸可能。
\begin{lstlisting}
int length = 0; // tail[] and a[] have capacity at least n.
for (int i = 0; i < n; i++) {
    int pos = lower_bound(tail, tail + length, a[i]) - tail;
    tail[pos] = a[i];
    if (pos == length) length++;
}
printf("%d\n", length);
\end{lstlisting}
时间 $O(n\log n)$。tail的内容不一定是一条真实的LIS；若要求恢复路径，还需要记录位置和前驱。不下降子序列对应upper\_bound，而不是lower\_bound。

若两个序列都是1到n的排列，把第一序列每个值的位置记为pos，再将第二序列映射成位置序列。共同子序列需要在第一序列位置上递增，因此转成LIS。这个转换利用了元素唯一的性质，不能直接用于有任意重复值的一般LCS。

\subsection{实数二分与三分的前提}
实数二分可固定迭代次数，避免直接比较浮点相等。例如连续函数在区间两端异号，可以保留仍有变号的半区间来寻找一个根。端点恰好为根时先处理；两个端点同号不代表没有根，也不能保证可以用这一规则寻找全部根。三次方程若有多个根，应先按题目给定性质隔离各根。

三分用于已知单峰或单谷的目标函数。对单峰求最大值，比较两个三等分点，舍弃不可能包含峰值的一侧；若函数存在多个峰，这个依据就不成立。分数规划、爬山和模拟退火还各有额外假设，不能用“都在搜索答案”把它们当作二分的替代模板。基础前三题训练应先把可证明的单调模型掌握牢固。

\subsection{可删堆的延迟删除}
优先队列通常只支持删除堆顶。若外部能保证删除请求对应一个当前存在的元素，可维护两个同类型堆：A保存所有插入值，B保存待删除值。每次取有效堆顶前，若两堆顶相等就同时弹出，直到不同或有堆为空。

重复值应按次数插入删除堆，不是用布尔标记；还要维护逻辑元素数，以区分真实为空与仍有待清理的节点。延迟删除让每个入堆元素最多被清理一次，但实现前应先判断multiset是否已经足够简单可靠。

\subsection{字符串匹配的线性预处理}
直接在每个起点比较模式串的复杂度最坏为 $O(nm)$，在小规模下完全可用。规模较大时，可用KMP的前缀函数维护“已匹配前缀还能退到多长”。设pi[i]为模式串前缀0到i的最长相等真前后缀长度：
\begin{lstlisting}
void prefix_function(const char s[], int pi[]) {
    int n = strlen(s);
    if (n == 0) return;
    pi[0] = 0;
    for (int i = 1; i < n; i++) {
        int j = pi[i - 1];
        while (j > 0 && s[i] != s[j]) j = pi[j - 1];
        if (s[i] == s[j]) j++;
        pi[i] = j;
    }
}
\end{lstlisting}
对非空模式p扫描文本时，用同样的回退规则更新匹配长度j；当j等于模式长度，记录一次匹配并回退到pi[j-1]，从而允许重叠匹配。由于指针前进与回退的总量受线性界约束，预处理和匹配合计 $O(n+m)$。空模式的匹配规则应另行规定，不能直接访问p[0]。

\subsection{附录练习与解答}
\textbf{练习1。}数量10如何拆包？\textbf{解答。}1、2、4、3，和为10，能表示0到10的每个件数。

\textbf{练习2。}两个排列 $[3,1,2]$ 与 $[1,3,2]$ 的LCS长度是多少？\textbf{解答。}第一排列位置为pos[3]=1、pos[1]=2、pos[2]=3；第二个映射为 $[2,1,3]$，LIS长度2，所以LCS长度2。

\textbf{练习3。}模式aaa在文本aaaa中出现几次？\textbf{解答。}2次，起点0和1。记录第一次后应保留可重叠的前后缀，而不是把匹配长度清成0。
\newpage
\section{资料来源与继续练习}
本讲义依据以下原始材料重新编排知识体系，补充了自包含教学例题、解答和参考程序。原课件中的重点概念问题已在正文采用修正后的表述，包括unique语义、DP状态信息、背包初始化、BFS归类和计数背包撤销。

\begin{longtable}{P{7.4cm}P{1cm}P{6cm}}
\toprule 原始文件 & 页数 & 本讲义对应内容\\\midrule\endhead
历年前三题考点梳理.pdf & 5 & 题型安排、基础模拟、综合模拟选题方向。\\
练习题与真题.pdf & 2 & 继续训练入口，仓库中另存链接索引。\\
stl.pptx & 42 & 第3章STL专章、第10章解析与可删堆。\\
前缀和、差分、二分、双指针.pptx & 34 & 第4--5章及搜索答案的选读内容。\\
搜索、剪枝、暴力.pptx & 13 & 第7章与模板依赖展开。\\
dp入门.pptx & 18 & 第8章、LIS与LCS扩展。\\
动态规划之背包问题.pptx & 53 & 第9章与背包提高专题。\\
\bottomrule
\end{longtable}

基础题进一步练习可选：排序 \prob{P1177}、去重 \prob{P1059}、前缀和 \prob{P8218}、差分 \prob{P2367}、二分 \prob{P2249}、切木条 \prob{P2440}、全排列 \prob{P1706}、马的遍历 \prob{P1443}、LIS \prob{B3637}、最大子段和 \prob{P1115}、01背包 \prob{P1048}、完全背包 \prob{P1616}、分组背包 \prob{P1757}。这些题可能有与教学题不同的输入、限制和输出要求，应重新阅读原题。

提高题可选：\prob{P2678}、\prob{P1219}、\prob{P3956}、\prob{P1776}、\prob{P1507}、\prob{P1064}、\prob{P1439}、\prob{P4141}。其中P4141的目标是删除指定物品后的恰好装满方案数末位数字，不是最大价值；P1439的较大数据需要利用排列性质，普通二维LCS不能覆盖全部规模。

接口语义可查阅C++工作草案的\href{https://eel.is/c++draft/sort}{sort条目}和\href{https://eel.is/c++draft/alg.unique}{unique条目}。平台原题和标准文档是进一步核对的依据，本讲义正文的教学例题与解答已经提供独立学习所需的规则和算法。

\subsection{团队题单与专题入口}
\textbf{练习总入口：}\href{https://www.luogu.com.cn/team/113242/training}{洛谷团队113242的训练题单}。以下5个专题地址来自原资料。团队页与专题题单目前需要访问权限，整理时未能核对其中的完整题目列表；正文中的公开题目链接不据此认定为团队题单成员。
\begin{longtable}{P{8cm}P{6cm}}
\toprule 专题 & 打开题单\\\midrule\endhead
二分、前缀和、差分、双指针 & \href{https://www.luogu.com.cn/training/1048062}{1048062}\\
搜索与剪枝 & \href{https://www.luogu.com.cn/training/1048981}{1048981}\\
STL & \href{https://www.luogu.com.cn/training/1055276}{1055276}\\
简单动态规划 & \href{https://www.luogu.com.cn/training/1052248}{1052248}\\
贪心 & \href{https://www.luogu.com.cn/training/1052246}{1052246}\\
\bottomrule\end{longtable}
\subsection{原资料真题做题入口}
以下按原资料中的第一题、第二题分组保留23个入口，便于直接进入平台作答。这是原资料链接索引；题名、认证编号及跨平台重复关系尚未全部核对，不将其伪标成已讲解的某道原题。正文的自编综合模拟也不等同于这些平台真题。
\begin{longtable}{P{1.5cm}P{2cm}P{10.5cm}}
\toprule 序号 & 分组 & 在线做题链接\\\midrule\endhead
1 & 第一题 & \url{https://www.acwing.com/problem/content/3250/}\\
2 & 第一题 & \url{https://www.acwing.com/problem/content/3260/}\\
3 & 第一题 & \url{https://www.acwing.com/problem/content/3275/}\\
4 & 第一题 & \url{https://www.acwing.com/problem/content/3290/}\\
5 & 第一题 & \url{https://www.acwing.com/problem/content/3295/}\\
6 & 第一题 & \url{https://www.acwing.com/problem/content/4009/}\\
7 & 第一题 & \url{https://www.acwing.com/problem/content/4702/}\\
8 & 第一题 & \url{https://www.acwing.com/problem/content/5019/}\\
9 & 第二题 & \url{https://www.acwing.com/problem/content/3266/}\\
10 & 第二题 & \url{https://www.acwing.com/problem/content/3286/}\\
11 & 第二题 & \url{https://www.acwing.com/problem/content/3296/}\\
12 & 第二题 & \url{https://www.acwing.com/problem/content/3415/}\\
13 & 第二题 & \url{https://www.acwing.com/problem/content/4010/}\\
14 & 第二题 & \url{https://www.acwing.com/problem/content/4458/}\\
15 & 第二题 & \url{https://www.acwing.com/problem/content/4513/}\\
16 & 第二题 & \url{https://www.acwing.com/problem/content/4703/}\\
17 & 第二题 & \url{https://www.acwing.com/problem/content/5020/}\\
18 & 第二题 & \url{https://www.acwing.com/problem/content/5723/}\\
19 & 第二题 & \url{https://oj.shumeng.tech/p/CSP202605B}\\
20 & 第二题 & \url{https://oj.shumeng.tech/p/CSP202409B}\\
21 & 第二题 & \url{https://oj.shumeng.tech/p/CSP202412B}\\
22 & 第二题 & \url{https://oj.shumeng.tech/p/CSP202503B}\\
23 & 第二题 & \url{https://oj.shumeng.tech/p/CSP202506B}\\
\bottomrule\end{longtable}
